\chapter{neuron 3.0 Architecture}\label{modern_architecture}
\index{neuron 3.0}

The 2026 reanimation preserves the object design described in this
Manifest while giving it three contemporary ways to be used: the
original menu interface, a local graphical interface, and a loopback
HTTP interface.  The same C++ model and statistical classes serve all
three.  Python tools prepare data and export trained models, but do not
duplicate the engine.

The central architectural rule is that every mechanism has one
authoritative implementation in the class that conceptually owns it.
The graphical interface may call that implementation, and a
cross-validation run may call it repeatedly, but neither copies its
mathematics.  Dependencies point down through the layers:

\begin{center}
\begin{tabular}{|>{\raggedright\arraybackslash}p{4.35in}|}
\hline
Human and programmatic interfaces: GUI, REST API, frozen CLI menus \\
\hline
Experiment orchestration: cross-validation, OBD sizing, reports,
automatic optimizer selection \\
\hline
Model and data interfaces: Model, DataSet, TwoSet, split plans,
model factory and clone helpers \\
\hline
Training and concrete methods: Iterative, Network, Logistic,
SimpleProp, BareProp, BackProp, LDFA, QDFA, RegressNet \\
\hline
Numerical foundation: Matrix, vector operations, statistics and
utility functions \\
\hline
\end{tabular}
\end{center}

A lower layer never learns a higher concern.  Matrix can gather rows
but does not know what a fold is.  DataSet can materialize a fold but
does not choose a model.  A model fits itself but does not run a
cross-validation loop.  OBD chooses a hidden-layer size but does not
own final evaluation.  Cross-validation owns repetition, not training.
This separation keeps the matrix notation, the implementation, and the
statistical report traceable to one another.

\section{Interfaces}

\subsection{Command line}

The interactive menus remain the authoritative feature inventory and
are intentionally frozen.  Reproducible command-line runs feed a
complete answer file to the engine and always provide a seed:

\begin{verbatim}
./build/neuron --seed 42 < session.in > session.out
\end{verbatim}

The transcript is the permanent record of the run.  The program is not
intended to be driven by keystrokes in an automated experiment.

\subsection{Graphical interface}

\begin{verbatim}
./build/neuron --gui
\end{verbatim}

This starts an embedded cpp-httplib server on
\texttt{127.0.0.1} at an operating-system-assigned port and opens a
single local web page.  The page provides the full load, model, train,
analysis, report, and save workflow.  It includes realtime training
progress and ROC plots.  The server is loopback-only: it is a local
application, not a network service.

GUI and CLI parity is a contract.  Every CLI capability must have both
a page control and an HTTP parameter.  The maintained option-by-option
map is \texttt{docs/gui\_cli\_parity.md}.  Features added beyond the
frozen CLI, including automatic optimizer selection, OBD sizing, and
cross-validation, are available through both the GUI and the HTTP
interface.

\subsection{Python tools}

\texttt{tools/mkdataset.py} turns raw CSV data into the numeric,
outcome-last format used by the engine and writes the key and
input-group files needed to interpret the encoded columns.
\texttt{tools/neuron2web.py} combines a saved network, scaling factors,
and a human-readable label specification into a self-contained HTML
calculator.  The deployed calculator performs inference entirely in
the browser.

\section{Data ownership and splitting}

\texttt{DataSet} owns loaded and normalized observations.
\texttt{TwoSet} owns classification and ROC reporting for a modeled
set.  The dedicated \texttt{split} layer owns index plans and provides:

\begin{itemize}
\item outcome-stratified train and test splits;
\item optional covariate stratification with quantile cells and
Hamilton apportionment;
\item group-aware splits that keep an entire cluster on one side;
\item three-way train, validation, and untouched test splits; and
\item outcome-stratified folds for cross-validation.
\end{itemize}

These mechanisms solve different problems.  Covariate stratification
makes two samples more alike and is useful for subgroup coverage.
Grouping makes evaluation harder by preventing site or cluster
leakage.  A validation set permits model selection without repeatedly
looking at the final test set.

\section{Models and training}

\texttt{Model} defines the common modeling interface.
\texttt{Network} adds connected layers and weights.
\texttt{Iterative} owns the training loop, stopping conditions,
progress callbacks, and convergence state.  Concrete model classes own
their distinct mathematics:

\begin{description}
\item[SimpleProp] One hidden layer with bias nodes.
\item[BareProp] One hidden layer without bias nodes.
\item[BackProp] Multiple hidden layers and multi-output networks.
\item[Logistic] Binary logistic regression, including Wald tests and
the design-matrix condition number.
\item[LDFA and QDFA] Linear and quadratic discriminant analysis.
\item[RegressNet] Forward and reverse stepwise variable selection
using the Wilks likelihood-ratio test.
\end{description}

Training supports canonical gradient descent, conjugate gradient
descent, and Shanno's method.  The automatic selector probes all three
from cloned starting weights for a short wall-clock budget, adopts the
lowest-error probe, and continues from the progress it made.

Stopping conditions are evaluated every iteration, independently of
how often an iteration row is printed.  Reaching the iteration ceiling
means that the requested operation ran, but the fit did not converge.
The weights remain available for continuation; an unfinished fit is
not eligible for model selection or cross-validation scoring.

\section{Model selection and evaluation}

\subsection{OBD hidden-layer sizing}

The \texttt{obd} layer grows a SimpleProp hidden layer while held-out
error improves, then prunes by saliency.  Each candidate size is
stopped at its held-out minimum or by a valid training stopping
condition.  With a three-way split, the search monitors validation and
leaves test untouched.  A successful standalone OBD run replaces the
current model with the selected, trained network.

\subsection{Cross-validation}

\texttt{crossval} applies logistic regression, LDFA, QDFA, a fixed
neural architecture, or nested OBD to one shared fold plan.
\texttt{cvadapters} connects those procedures to their existing model
implementations.  \texttt{cvreport} renders a three-tier result:

\begin{enumerate}
\item a headline comparison suitable for decisions;
\item a detailed human-readable audit; and
\item machine-readable prediction, metric, and run files.
\end{enumerate}

Nested OBD repeats architecture selection inside each training fold.
It therefore estimates the performance of the selection procedure,
not the performance of an architecture chosen using all observations.
Cross-validation is a standalone analysis and does not replace the
current model.

\subsection{Locked-test inference}

A locked test set may be removed before cross-validation.  Procedures
are compared on the development rows, refitted by their own rules, and
scored once on the locked rows.  \texttt{delong} supplies paired AUC
confidence intervals and the prespecified AUC contrast when the caller
declares independent rows.  Clustered observations require
cluster-aware inference, so ordinary DeLong inference is withheld
rather than silently reported.

\section{Statistical reporting}

Binary models report classification accuracy, sensitivity,
specificity, and two complementary ROC estimates.  The trapezoidal
area is the exact non-parametric AUC over every distinct operating
point.  The primary binormal $A_z$ is fitted to the z-ROC curve and is
reported with a stratified bootstrap confidence interval.  Reports
state the number of fitted operating points, bootstrap resamples, and
failures.  Logistic models additionally report coefficient-level Wald
tests and the condition number.

\section{Run artifacts and audit trail}

The directory of the first loaded dataset is the run directory.
Networks, scaling factors, normalized sets, guesses, reports,
\texttt{neuron.log}, and \texttt{neuron\_actions.log} are written
there.  The latter records every GUI or API action with a timestamp and
its parameter values.  Long operations publish structured progress
through one shared status channel and honor graceful cancellation at
an iteration boundary.

\section{Source map for the modern layers}

\begin{longtable}{|>{\raggedright\arraybackslash}p{1.05in}|>{\raggedright\arraybackslash}p{3.15in}|}
\hline
\textbf{Files} & \textbf{Responsibility} \\ \hline
\endfirsthead
\hline
\textbf{Files} & \textbf{Responsibility} \\ \hline
\endhead
\texttt{gui.*},
\texttt{gui\_page.html} &
Local server, REST handlers, asynchronous job state, and browser
interface. \\ \hline
\texttt{split.*} &
Stratified, group-aware, three-way, and k-fold index plans. \\ \hline
\texttt{crossval.*},
\texttt{cvadapters.*},
\texttt{cvreport.*} &
Procedure comparison, connections to model implementations, and
human plus machine-readable reports. \\ \hline
\texttt{obd.*} &
Grow-and-prune hidden-layer selection. \\ \hline
\texttt{autoalgo.*},
\texttt{plateau.h} &
Optimizer probing and plateau detection shared by training
workflows. \\ \hline
\texttt{delong.*} &
Paired locked-test AUC inference. \\ \hline
\texttt{modelfactory.*},
\texttt{netclone.*} &
Construction and faithful cloning through the common model
interfaces. \\ \hline
\texttt{dataset.*},
\texttt{twoset.*} &
Observation ownership, normalization, classification, ROC, and
statistical reporting. \\ \hline
\texttt{model.*},
\texttt{network.*},
\texttt{iterative.*} &
Core model contract, network structure, and authoritative training
loop. \\ \hline
\texttt{matrix.*},
\texttt{vector\_ops.*},
\texttt{stats.*} &
Bounds-checked numerical and statistical foundation. \\ \hline
\end{longtable}
